\section{Discussion}\label{discussion}
This paper develops two extensions of the estimating-equation variance
framework of \citet{robins2000inference} to settings that arise
routinely in modern medical data
analysis: chained conditional (MICE) imputation and predictive mean
matching.  The extensions are not ad hoc adjustments but follow directly
from the RW construction once the relevant score and nuisance objects are
identified for each setting.  Both extensions are validated in simulation
and the key implementation objects are exposed by a companion \textsf{R}
package, \texttt{rw}, that seamlessly integrates with \texttt{mice}
package output.  Our open-source software is easy to use because it
mirrors the functions already in place for Rubin's rules variance
estimators and takes as input a \texttt{mice} package imputed dataset.
The failure of Rubin's rules under uncongeniality has been understood
since \citet{meng1994multiple} and formalised for estimating equation
analyses by \citet{wang1998large}.  The present work bridges these
theoretical results and applied practice by showing that both failure
modes (anti-conservative and conservative) arise in common MICE
and PMM workflows and can be corrected within the same sandwich
framework.  The GS validation setting, where eligibility is determined
by the imputed variable and Rubin's rules overstate the variance by
approximately 30--50\%, is an instance of Meng's conservative
uncongeniality that has become increasingly common in EHR based analyses
with imputed inclusion criteria \citep{shen2026missingness}.  The
RW-MICE and RW-PMM estimators restore calibration in this setting
without requiring any change to the imputation or analysis workflow.
\paragraph{The bootstrap.}
As noted in Section~\ref{introduction}, the bootstrap is an alternative
to analytic variance estimation under uncongeniality: re-running the
full imputation and analysis procedure across bootstrap samples
consistently estimates the variance of the implemented estimator
regardless of congeniality
\citep{schomaker2018bootstrap, bartlett2020bootstrap}.  The RW estimator
has a practical advantage when the imputation workflow is expensive to
repeat.
\paragraph{Limitations.}
Several limitations should be noted.  First, the PMM extension requires
donor indices \(J_i^{(p)}\) to be recorded during imputation; they
cannot be reconstructed from completed datasets.  This is a software
requirement rather than a theoretical obstacle, and the companion package
handles it automatically, but practitioners using standard
\texttt{mice} output will need to update their imputation workflow.
Second, the present work considers the case where one continuous
variable is imputed by the smooth copied donor law and the remaining
blocks are parametric.  When several variables are imputed by PMM with
potentially different donor maps, the multi-source covariance structure
requires additional work.  The analysis score contributions of two
recipients can be correlated through two independent donor maps
simultaneously, and stacking the donor source representations from
different PMM blocks does not automatically account for this joint
dependence.  Possible extensions include block PMM with common donors or
a multi-way source aggregation; these fit within the same
source-aggregation framework, and we plan to support general PMM
problems in a future release of the package.
Third, the smooth copied donor law is a working device for computing the
RW score; it does not change the imputed values.  The asymptotic
equivalence between the RW estimator for the smooth law and the
bootstrap variance of the hard nearest-neighbor PMM estimator has not
been formally established.  The simulation evidence is supportive, but a
rigorous analysis of the approximation error as \(k\to\infty\) or
\(\lambda\to\infty\) remains an open problem.
Fourth, the paper focuses on Gaussian and logistic analysis estimating
equations.  Extensions to survival analysis, competing risks, or
semiparametric models with partially observed outcomes are not
considered, though the RW framework is in principle applicable to any
regular estimating equation.